\item \textbf{Índices de refracción de tres láminas de plástico.}
Analice la siguiente situación en su grupo: Tres láminas de plástico tienen índices de refracción desconocidos. Se colocan una sobre la otra de dos en dos y se dirige un rayo láser a las hojas desde arriba. En los tres casos, el rayo láser se ajusta de modo que golpee la interfaz entre la lámina superior y la inferior (no entre el aire y la lámina superior) en un ángulo de $26.5^\circ$ con la normal.
\begin{enumerate}[label=(\roman*)]
    \item La hoja 1 se coloca encima de la hoja 2. El haz refractado dentro de la hoja 2 forma un ángulo de $31.7^\circ$ con la normal.
    \item La lámina 3 se coloca encima de la lámina 2. El haz refractado dentro de la lámina 2 forma un ángulo de $36.7^\circ$ con la normal.
    \item La lámina 1 se coloca encima de la lámina 3. El haz refractado dentro de la lámina 3 forma un ángulo de $23.1^\circ$ con la normal.
\end{enumerate}
Ahora genere una respuesta de consenso en su grupo: ¿Tiene suficiente información para encontrar el índice de refracción para cada hoja?